% Critical-weight estimates for the original-norm class. Input from sec_printed.tex; macros are prefixed \os.
\newcommand{\osomc}{\omega_c}        % critical weight X_0^{1/2} l^c
\newcommand{\osomp}{\omega}          % periodization weight X_0^{1/2} l^{c'}
\newcommand{\osX}{X_0}               % X_0 = cosh d(x_0,\cdot)
\newcommand{\osL}{\ell}              % l = log(2X_0)/log 2
\newcommand{\osHK}{\widetilde p}     % continuation of the heat kernel to complex distance
\newcommand{\osLam}{\Lambda}         % Lambda = (-Delta)^{1/2}

\subsection{Critical-weight multipliers}\label{app:o1st}

To handle the original-norm class we work at the weight $X_0^{1/2}$ with a logarithmic correction, in the notation of
Section~\ref{subsec:pc-weights}. A multiplier theorem at this weight (Theorem~\ref{os:cw}) gives a factorization
$a=\ue^{-\sigma\Lambda_m}\tilde f$ with critically weighted data $\tilde f$ (Proposition~\ref{os:data}); word bounds
(Lemma~\ref{os:words}) and transfer over shells (Section~\ref{subsec:pc-transfer}) then control the derivatives, flow and defect
appearing in the remainder.

The multiplier theorem rests on a balance between two exponential factors. Moving mass from a shell $A_n$ to a shell $A_k$
with $k>n$ costs the factor $\ue^{(k-n)/2}$ of the weight $X_0^{1/2}$. By finite propagation speed, only the part of the
multiplier with $|r|\ge k-n-1$ contributes, and its tail supplies the factor $\ue^{-(k-n)/2}(1+k-n)^{-3/2}$. What remains is a
polynomially decaying kernel in $k-n$, and the logarithmic exponent $\tfrac12<c<1$ makes it summable. The profiles in
Appendix~\ref{app:profiles} verify the tail condition for the Poisson semigroup and for the change of spectral parameter.

\begin{lemma}[finite propagation speed]\label{os:speed}
Let $f,h\in L^2(\Hb)$ with $\operatorname{supp}f\subset\{\tilde\psi\le R_1\}$ and $\operatorname{supp}h\subset\{\tilde\psi\ge R_2\}$. Then
$\langle h,\cos(rD)f\rangle=0$ and $\langle f,\cos(rD)h\rangle=0$ for $|r|<R_2-R_1$.
\end{lemma}
\begin{proof}
For $\lambda>0$ apply Lemma~\ref{lem:wave}(b) to $W_{\lambda,n}$, with $l=\lambda$, and let $n\to\infty$. This gives
\[
\norm{\osX^\lambda\cos(rD)f}\le\ue^{(\lambda+1/2)|r|}\norm{\osX^\lambda f}\le\ue^{(\lambda+1/2)|r|+\lambda R_1}\norm f.
\]
Hence
$\norm{1_{\{\tilde\psi\ge R_2\}}\cos(rD)f}\le\ue^{|r|/2}\ue^{-\lambda(R_2-R_1-|r|)}\norm f$, which tends to $0$ as $\lambda\to\infty$. The second identity follows
because $\cos(rD)$ is self-adjoint.
\end{proof}

\begin{theorem}[multiplier theorem at the critical weight]\label{os:cw}
Let $\tfrac12<c<1$. Let $g\in L^1(\mathbb R)$ be real with
\[
\int_{|r|\ge T}|g(r)|\,dr\le B\,\ue^{-T/2}(1+T)^{-3/2}\qquad\text{for all }T\ge0,
\]
and let $\mathsf G:=\int_{\mathbb R}g(r)\cos(rD)\,dr$, a bounded operator on $L^2(\Hb)$. Then for all $f\in L^2(\Hb)$,
\[
\norm{\osomc\,\mathsf Gf}\le C_c\,B\,\norm{\osomc f},
\]
with $C_c$ depending only on $c$.
\end{theorem}
\begin{proof}
Put $f_n=1_nf$ and $a_n=\norm{\osomc f_n}$, and let $T(k,n):=\max(|k-n|-1,0)$. By Lemma~\ref{os:speed}, for $h$ supported in $A_k$ the function
$r\mapsto\langle h,\cos(rD)f_n\rangle$ vanishes on $|r|<T(k,n)$. Hence
$\norm{1_k\mathsf Gf_n}\le\int_{|r|\ge T(k,n)}|g|\,\norm{f_n}\le B\,\ue^{-T/2}(1+T)^{-3/2}\norm{f_n}$, and by \eqref{os:shell}
\[
\begin{gathered}
\norm{\osomc1_k\mathsf Gf_n}\le\mathsf B(k,n)\,a_n,\\
\mathsf B(k,n):=C\,B\,\Bigl(\frac{1+k}{1+n}\Bigr)^c(1+|k-n|)^{-3/2}\times
\begin{cases}1,&k\ge n,\\ \ue^{-(n-k)},&k<n.\end{cases}
\end{gathered}
\]
Since $\mathsf Gf=\sum_n\mathsf Gf_n$ in $L^2$, $\norm{\osomc\mathsf Gf}^2\le\sum_k\bigl(\sum_n\mathsf B(k,n)a_n\bigr)^2$. We apply the Schur test on
$\mathbb N_0$ with the weights $p_n=(1+n)^{-1/2}$.
\begin{itemize}
\item $\sum_n\mathsf B(k,n)p_n\le C'B\,p_k$. For $n\ge k/2$: $((1+k)/(1+n))^c\le2^c$, $p_n\le\sqrt2\,p_k$, and $\sum_j(1+j)^{-3/2}<\infty$. For $n<k/2$:
$\mathsf B(k,n)p_n\le CB(1+k)^{c-3/2}(1+n)^{-c-1/2}$; the sum over $n$ converges because $c>\tfrac12$, and $(1+k)^{c-3/2}\le p_k$ because $c<1$.
\item $\sum_k\mathsf B(k,n)p_k\le C'B\,p_n$. For $k<n$: $((1+k)/(1+n))^cp_k\le p_n$ because $c>\tfrac12$. For $n\le k\le2n+1$: $((1+k)/(1+n))^c\le2^c$ and
$p_k\le p_n$. For $k>2n+1$: $\mathsf B(k,n)p_k\le CB(1+n)^{-c}(1+k)^{c-2}$, and $\sum_{k>2n+1}(1+k)^{c-2}\le C(1+n)^{c-1}$ because $c<1$.
\end{itemize}
So $\sum_k(\sum_n\mathsf B(k,n)a_n)^2\le(C'B)^2\sum_na_n^2=(C'B)^2\norm{\osomc f}^2$.
\end{proof}

The cosine-transform representations and tail bounds needed to apply Theorem~\ref{os:cw}
are given in Lemma~\ref{os:prof}, proved in Appendix~\ref{app:profiles}.
They give the following two operator bounds: Poisson smoothing preserves the critical weight,
and the change from $\Lambda$ to $\Lambda_m$ is bounded at that weight.

\begin{corollary}[Poisson smoothing and the shifted multiplier at the critical weight]\label{os:cwp}
Let $\tfrac12<c<1$. For $0<\tau\le1$ and $f\in L^2$, $\norm{\osomc\,\ue^{-\tau\osLam}f}\le k_c\norm{\osomc f}$. For $m\ge\tfrac12$,
$\norm{\osomc(\Lambda_m-\osLam)f}\le k_{m,c}\norm{\osomc f}$.
\end{corollary}
\begin{proof}
Theorem~\ref{os:cw} with $g=g^{1/2}_\tau$ and Lemma~\ref{os:prof}(i),(iv), respectively with $g=G_m$ and Lemma~\ref{os:prof}(v).
\end{proof}

\subsection{Factorization and derivative estimates for the data}\label{subsec:pc-data}
We now factor the Poisson-smoothed amplitude through the shifted semigroup and bound its derivatives at a power weight.
The next subsection transfers these bounds to the critical weight used for periodization.

\begin{proposition}[factorization of the data]\label{os:data}
Let $\tfrac12<c<1$, $m\ge\tfrac12$ and $0<\sigma<\alpha\le1$ with $\sigma m\le\tfrac12$, and put
\[
\mathsf R:=\ue^{-(\alpha-\sigma)\osLam}\,\ue^{\sigma(\Lambda_m-\osLam)},
\]
a bounded operator. Then $\ue^{-\sigma\Lambda_m}\mathsf R=\ue^{-\alpha\osLam}$, and
$\norm{\osomc\mathsf Rf}\le C_{\mathrm d}\norm{\osomc f}$ for all $f\in L^2$, with a constant $C_{\mathrm d}$ depending only on $m$ and $c$.
Consequently, if $\beta>\tfrac12$ and $\pcNorm a<\infty$ at $(\alpha,\beta)$, then $\tilde f:=\mathsf R\,\ue^{\alpha\osLam}a$ satisfies
\[
a=\ue^{-\sigma\Lambda_m}\tilde f,\qquad\norm{\osomc\tilde f}\le C\,\pcNorm a,
\]
with $C=C(m,c,\beta)$.
\end{proposition}
\begin{proof}
The identity is the bounded functional calculus: $-\sigma\Lambda_m+\sigma(\Lambda_m-\osLam)-(\alpha-\sigma)\osLam=-\alpha\osLam$. The operator
$\mathsf B:=\Lambda_m-\osLam$ is bounded, and $\ue^{\sigma\mathsf B}=\sum_k\sigma^k\mathsf B^k/k!$ in operator norm. By Corollary~\ref{os:cwp},
$\norm{\osomc\mathsf B^kf}\le k_{m,c}^k\norm{\osomc f}$, so the partial sums are bounded in the weighted norm by $\ue^{\sigma k_{m,c}}\norm{\osomc f}$, and
Fatou's lemma gives $\norm{\osomc\ue^{\sigma\mathsf B}f}\le\ue^{k_{m,c}}\norm{\osomc f}$. Then apply Corollary~\ref{os:cwp} to
$\ue^{-(\alpha-\sigma)\osLam}$; so $C_{\mathrm d}=k_c\,\ue^{k_{m,c}}$ works. For the consequence put $f=\ue^{\alpha\osLam}a$, so $a=\ue^{-\alpha\osLam}f$. Since $\tilde\psi\le d\le\langle d\rangle+1$ and
$\beta>\tfrac12$, $\osomc\le\ue^{d/2}(1+d/\log2)\le C_\beta\,\ue^{\beta\langle d\rangle}$.
\end{proof}

For the derivative bounds write $c_H:=1+4/\ue$, the constant in the profile bound of Lemma~\ref{os:prof}(iii).

\begin{lemma}[words of the shifted Poisson semigroup at a bounded weight]\label{os:words}
Let $m\ge\tfrac12$ and $\rho\ge0$ with $2(\rho+\tfrac12)^2<m^2$, and let $W$ be a bounded weight (Section~\ref{app:ws}) with
$(X\log W)^2+(Y\log W)^2\le\rho^2$. Put $\mu:=m^2-\tfrac14$, $C_6:=2\ue\,c_H$ and $Q:=4+2\rho$.
\begin{itemize}
\item[(a)] For $\tau,s>0$ and $u\in L^2$:
\[
\norm{W\ue^{-\tau\Lambda_m}u}\le\sqrt2\norm{Wu},\qquad
\norm{W\Lambda_m\ue^{-s\Lambda_m}u}\le(2c_H/s)\norm{Wu}.
\]
Hence for $j\ge1$ and $\sigma>0$, $\norm{W(\mu-\Delta)^j\ue^{-\sigma\Lambda_m}u}\le(4jc_H/\sigma)^{2j}\norm{Wu}$.
\item[(b)] For $0<\sigma\le1$, every word $Z_S$ of length $k$ in $X,Y$, and $h\in L^2$,
\[
\norm{WZ_S\,\ue^{-\sigma\Lambda_m}h}\le\sqrt2\,Q^k\,k!\,(C_6/\sigma)^k\,\norm{Wh}.
\]
\end{itemize}
In particular, let $\lambda\ge0$ with $2(\lambda+\tfrac12)^2<m^2$ and $\osX^\lambda h\in L^2$. Then, with $\rho=\lambda$ in the definition of $Q$,
\[
\norm{\osX^\lambda Z_S\,\ue^{-\sigma\Lambda_m}h}\le\sqrt2\,Q^k\,k!\,(C_6/\sigma)^k\,\norm{\osX^\lambda h}.
\]
Moreover $\norm{\ue^{\sigma\osLam}\ue^{-\sigma\Lambda_m}}\le1$.
\end{lemma}
\begin{proof}
(a) By Lemma~\ref{os:prof}(i), Lemma~\ref{lem:wave}(b) and Lemma~\ref{os:prof}(ii) with $\rho+\tfrac12$ in the role of $\lambda$,
$\norm{W\ue^{-\tau\Lambda_m}u}\le\int g^m_\tau(r)\ue^{\lambda|r|}dr\,\norm{Wu}\le\sqrt2\norm{Wu}$. The same with $h^m_s$ and Lemma~\ref{os:prof}(iii) gives the second
bound. For the third write $(\mu-\Delta)^j\ue^{-\sigma\Lambda_m}=(\Lambda_m\ue^{-(\sigma/2j)\Lambda_m})^{2j}$.

(b) Let $v_j:=(\mu-\Delta)^j\ue^{-\sigma\Lambda_m}h$. Since $\Lambda_m\ge\osLam$, $v_j\in\operatorname{Dom}(\ue^{\sigma'\osLam})$ for $\sigma'<\sigma$, so $v_j$ is smooth
and all words applied to $v_j$ lie in $L^2$ (Lemmas~\ref{lem:wc} and~\ref{lem:lift}). Put
\[
A_{k,j}:=\max_{|S|=k}\norm{WZ_Sv_j},\qquad\mathsf B_{k+1,j}:=\max_{|S|=k}\bigl(\norm{WXZ_Sv_j}^2+\norm{WYZ_Sv_j}^2\bigr)^{1/2},
\]
so $A_{k+1,j}\le\mathsf B_{k+1,j}$. Fix $S$ with $|S|=k$ and $w=Z_Sv_j$. Integration by parts gives
\begin{multline*}
\norm{WXw}^2+\norm{WYw}^2=\mathrm{Re}\langle W^2w,(\mu-\Delta)w\rangle-\mu\norm{Ww}^2\\
-\mathrm{Re}\langle X(W^2)w,Xw\rangle-\mathrm{Re}\langle Y(W^2)w,Yw\rangle .
\end{multline*}
Since $[\Delta,X]=2XY$ and $[\Delta,Y]=-2X^2$, $[\Delta,Z_S]$ is a sum of $k$ words of length $k+1$ with coefficients $\pm2$, and
$(\mu-\Delta)w=Z_Sv_{j+1}-[\Delta,Z_S]v_j$. The last two terms are at most $2\rho\norm{Ww}(\norm{WXw}^2+\norm{WYw}^2)^{1/2}$ together. Hence
\[
\mathsf B_{k+1,j}^2\le A_{k,j}A_{k,j+1}+(2k+2\rho)A_{k,j}\,\mathsf B_{k+1,j},
\]
and therefore
\[
A_{k+1,j}\le\mathsf B_{k+1,j}\le(2k+2\rho)A_{k,j}+(A_{k,j}A_{k,j+1})^{1/2}.
\]
We claim $A_{k,j}\le\sqrt2\,Q^k(C_6/\sigma)^{k+2j}(k+2j)!\,\norm{Wh}$ for all $k,j\ge0$. For $k=0$ this is (a), because
$(4jc_H/\sigma)^{2j}\le(2\ue c_H/\sigma)^{2j}(2j)!$. If it holds for $k$ and all $j$, then with $n=k+2j$, using $(n!\,(n+2)!)^{1/2}\le\sqrt2\,(n+1)!$,
$(2k+2\rho)\,n!\le(2+2\rho)(n+1)!$ and $(C_6/\sigma)^n\le(C_6/\sigma)^{n+1}/C_6$,
\[
A_{k+1,j}\le\sqrt2\,Q^k(C_6/\sigma)^{n+1}(n+1)!\,\bigl((2+2\rho)/C_6+\sqrt2\bigr)\norm{Wh},
\]
which is the claim for $k+1$, because $C_6\ge1$ and so $(2+2\rho)/C_6+\sqrt2\le Q$. The case $j=0$ is (b). The statement for $\osX^\lambda$ follows with $W=W_{\lambda,n}$ and $n\to\infty$, and the last
statement from $\Lambda_m\ge\osLam$.
\end{proof}

\subsection{Transfer and periodization at the critical weight}\label{subsec:pc-transfer}
The transfer estimate below applies to the derivative, flow and defect maps in the approximation proof.
Periodization at the same weight then turns these estimates into bounds on the quotient.

\begin{lemma}[transfer over shells]\label{os:transfer}
Let $0\le c\le1$ and $C_{\mathrm{sh}}>0$, and let $\varpi$ be a measurable weight with $\ue^{n/2}(1+n)^c\le\varpi\le C_{\mathrm{sh}}\,\ue^{n/2}(1+n)^c$ on $A_n$ for
every $n\ge0$; by \eqref{os:shell}, $\varpi=\osomc$ with $C_{\mathrm{sh}}=3\ue^{1/2}$ is an example. Let $T:L^2(\Hb)\to L^2(\Hb)$ be additive, and let
$0<M_0\le M_1$ be such that
\begin{itemize}
\item[(a)] $\norm{Tu}\le M_0\norm u$ for all $u\in L^2$, and
\item[(b)] $\norm{\osX\,Tu}\le M_1\norm{\osX u}$ for every $u$ supported in one shell $A_n$.
\end{itemize}
Then for every $h\in L^2$,
\[
\norm{\varpi\,Th}\le C_{\mathrm{tr}}\,M_0^{1/4}M_1^{3/4}\,\norm{\varpi h},
\]
with a constant $C_{\mathrm{tr}}$ depending only on $c$ and $C_{\mathrm{sh}}$.
\end{lemma}
\begin{proof}
Let $u=1_nh$. By (a), $\norm{1_kTu}\le M_0\norm u$. By (b), since $\osX\ge\ue^k$ on $A_k$ and $\osX\le\ue^{n+1}$ on $A_n$,
$\norm{1_kTu}\le\ue\,M_1\ue^{-(k-n)}\norm u$. For $k\ge n$, $\min(x,y)\le x^{1/4}y^{3/4}$ gives
$\norm{1_kT1_n}\le\ue^{3/4}M_0^{1/4}M_1^{3/4}\ue^{-3(k-n)/4}$; for $k<n$ we use $M_0\le M_0^{1/4}M_1^{3/4}$. By the bounds for $\varpi$ and $c\le1$,
$\sup_{A_k}\varpi/\inf_{A_n}\varpi\le C_{\mathrm{sh}}\,\ue^{(k-n)/2}(1+(k-n)_+)$. So with $a_n:=\norm{\varpi1_nh}$,
\[
\norm{\varpi1_kT1_nh}\le M_0^{1/4}M_1^{3/4}\,w_{k-n}\,a_n,\qquad w_j:=\begin{cases}C_{\mathrm{sh}}\ue^{3/4}(1+j)\ue^{-j/4},&j\ge0,\\C_{\mathrm{sh}}\ue^{-|j|/2},&j<0.\end{cases}
\]
Since $T$ is additive and bounded, $Th=\sum_nT1_nh$ in $L^2$, and therefore
$\norm{\varpi Th}^2\le\sum_k\bigl(\sum_n\norm{\varpi1_kT1_nh}\bigr)^2\le M_0^{1/2}M_1^{3/2}\bigl(\sum_jw_j\bigr)^2\sum_na_n^2$ by Young's inequality for
sequences.
\end{proof}

\begin{proposition}[periodization and transport at the critical weight]\label{os:p1}
Let $\Gamma$ be any Fuchsian group, $\tfrac12<c'\le1$, $\osomp=\omega_{c'}$, $t\ge0$, $\hbar>0$, $\varphi=\varphi_\infty=-\log y$, and
$w\in L^2(\Hb)$ with $\osomp w\in L^2(\Hb)$. Then the sum defining $[(S(t)w)\ue^{i\varphi/\hbar}]_\Gamma$ converges absolutely almost everywhere, and
\[
\bigl\|[(S(t)w)\ue^{i\varphi/\hbar}]_\Gamma\bigr\|_{L^2(M)}\le C(\Gamma,c')\,\ue^{t/2}(1+t/\log2)^{c'}\,\|\osomp w\|_{L^2(\Hb)}.
\]
In particular $\osomp$ satisfies \textup{(W1)} and \textup{(W2)} of Section~\ref{subsec:dd-surf}, with
$\osomp(z)/\osomp(\delta_sz)\le\ue^{|s|/2}(1+|s|/\log2)^{c'}\le(2/\log2)\,\ue^{|s|}$.
\end{proposition}
\begin{proof}
The proof is that of Lemma~\ref{lem:t2}, with $\osomp$ in place of $\ue^{s\tilde\psi}$; only the first and last steps change.

\emph{Pointwise Cauchy--Schwarz.} Since $\osomp(z')^2\asymp\ue^{d(z',x_0)}(1+d(z',x_0))^{2c'}$ for $z'\in\Hb$ and $d(\gamma z,x_0)=d(z,\gamma^{-1}x_0)$,
reindexing $\gamma\mapsto\gamma^{-1}$ and the counting bound of Lemma~\ref{lem:count} for the orbit $\Gamma x_0$ give
\[
\sum_\gamma\osomp(\gamma z)^{-2}\le C\sum_\gamma\ue^{-d(z,\gamma x_0)}(1+d(z,\gamma x_0))^{-2c'}\le C\sum_{n\ge0}\ue^{n}\ue^{-n}(1+n)^{-2c'}<\infty,
\]
uniformly in $z$, since $2c'>1$. Cauchy--Schwarz and unfolding, as in Lemma~\ref{lem:t2}, give
$\|[f]_\Gamma\|^2_{L^2(M)}\le C\int_\Hb\osomp^2|f|^2\,d\nu$ for $f=(S(t)w)\ue^{i\varphi/\hbar}$, with almost everywhere absolute convergence.

\emph{Transport.} $(S(t)w)(z)=\ue^{-t/2}w(\delta_tz)$, and $d(z,\delta_sz)=|s|$. Since $\tilde\psi$ is $1$-Lipschitz, $\osX(z)\le\ue^{|s|}\osX(\delta_sz)$ and
$\osL(z)\le\osL(\delta_sz)+|s|/\log2\le\osL(\delta_sz)(1+|s|/\log2)$; this is the ratio bound. Since $S(t)$ is unitary,
$\int\osomp^2|S(t)w|^2d\nu\le\ue^{t}(1+t/\log2)^{2c'}\int\osomp^2|w|^2d\nu$. Combining gives the claim.
\end{proof}

\begin{lemma}[the quotient propagator does not depend on the weight]\label{os:indep}
Let $\varpi_1,\varpi_2$ be two weights satisfying \textup{(W1)}, \textup{(W2)}. For $g\in L^2_{\varpi_1}\cap L^2_{\varpi_2}$ the element
\[
 \mathcal V_sg=\ue^{-is/(2\hbar)}\cU(-s)[\ue^{i\varphi/\hbar}S(s)g]_\Gamma\in L^2(M)
\]
is the same for both weights. In particular this applies to $\varpi_1=\osX^{\lambda}$ with $\lambda>\tfrac12$ and
$\varpi_2=\osomp$, and $L^2_{\osX^\lambda}\subset L^2_{\osomp}$.
\end{lemma}
\begin{proof}
For both weights the periodization is the same almost everywhere absolutely convergent sum, by \textup{(W1)}.
The inclusion holds because $\osomp=\osX^{1/2}\osL^{\,c'}\le C_\lambda\osX^\lambda$ for $\lambda>\tfrac12$, since $\osL=1+\tilde\psi/\log2$
grows like $\log\osX$.
\end{proof}
